\section{置信度属于程序，而非这次参数}
\label{sec:11-Mathematical-Statistics/04-Interval-Estimation/01}

你跑完实验，得到 95\% 置信区间 \([0.86,\,0.97]\)。写报告时你写下：``参数 \(\theta\) 有 95\% 的概率落在这个区间内。''

这个说法在直觉上毫无问题，在日常交流里人人能懂。但在频率学派的逻辑里，它在技术上是错的。而且这个对错不是咬文嚼字---它会直接影响你对``这个区间到底意味着什么''的理解。

\subsection{概率归属于区间，不归属于参数}

频率学派置信区间是一套程序：每当你拿到新的随机样本 \(X\)，程序吐出一个随机区间

\[
C(X)=[L(X),\,U(X)].
\]

这个程序的长期属性是：对每一个可能的参数值 \(\theta\)，

\[
P_\theta\{\theta\in C(X)\}\ge 1-\alpha.
\]

注意概率写在 \(P_\theta\) 下，作用于 \(C(X)\)。\(C(X)\) 是随机的，因为 \(X\) 是随机的。\(\theta\) 是固定的。长期来看，如果你无数次地从同一个 \(\theta\) 下抽样本、算区间，大约 \(1-\alpha\) 比例的区间会盖住那个固定的 \(\theta\)。

现在你把一次具体的实验做完了。数据变成了具体的数字，区间变成了具体的端点 \([0.86,0.97]\)。在频率学派的模型里，\(\theta\) 仍然是一个固定但未知的数。事件``\(\theta\) 是否在 \([0.86,0.97]\) 里''已经非真即假，只是你不知道答案。它不再拥有由抽样随机性定义的 95\% 概率---因为随机性已经用完了。

你可以打一个渔网的比方。撒网之前，你说``这张网有 95\% 的概率能捞到鱼''，这是一个关于程序的陈述。但网已经撒下、捞起来了，里面要么有鱼要么没有。你只是还没打开看。你不能说``网里的鱼此刻有 95\% 的概率存在''。

所以频率学派更准确的陈述是：

\begin{quote}
我们使用了一个在假设模型下长期覆盖率为 95\% 的程序；在这一次数据上，它给出了区间 \([0.86,\,0.97]\)。
\end{quote}

日常报告把它简写成``95\% 置信区间''，但在自己心里需要把它和``参数有 95\% 概率在其中''区分开。后者是 Bayesian 后验区间的逻辑，需要在参数上放置概率分布。频率学派没有给参数配概率，所以做不出那个陈述。

\subsection{Wald 区间：简单到危险}

把标准误里的未知参数直接替换成估计值，是最顺手也最容易出错的做法。对 Bernoulli 样本比例 \(\hat\theta=S/n\)，中心极限定理给出 \(\hat\theta\) 的渐近正态性。把 \(\sqrt{\theta(1-\theta)/n}\) 里的 \(\theta\) 替换为 \(\hat\theta\)（plug-in 近似），得到 Wald 区间：

\[
\hat\theta \pm z_{1-\alpha/2}\sqrt{\frac{\hat\theta(1-\hat\theta)}{n}},
\]

其中 \(z_{1-\alpha/2}\) 是标准正态的 \(1-\alpha/2\) 分位数，即 \(\Phi(z_{1-\alpha/2})=1-\alpha/2\)。

这个区间极其好算。但它有三重失败。

第一重：边界爆炸。取 \(n=20\)，假设你一次正面也没看到，\(S=0\)，\(\hat\theta=0\)。标准误估计值是 \(\sqrt{0\cdot 1/20}=0\)。你的 95\% 置信区间变成了 \([0,0]\)---一个长度为零的区间，宣称对覆盖 \(\theta\) 有 95\% 信心。这显然荒谬。零次成功只告诉你概率可能较小，不意味着它精确为零。这条区间对任何小的正 \(\theta>0\) 而言都完全没有覆盖，远不到标称的 95\%。

第二重：超出合法范围。\(\hat\theta\) 靠近边界时，加减一项可能把区间端点推出 \([0,1]\) 之外。你可以手工截断，但截断后的区间覆盖率不等于原标称值。

第三重：即便 \(\hat\theta\) 不极端，实际覆盖率也常低于标称值，尤其在 \(n\) 不够大或 \(\theta\) 靠近 0 或 1 时。所谓的``95\%''经常只是一个渐近保证，有限样本下可能是 80\% 或更低。

\subsection{不止 Wald 一种选择}

更好的 Bernoulli 比例区间包括：

Wilson score 区间：从 score 检验反解，不靠 plug-in 近似标准误。它不会退化成零长度，端点自动在 \([0,1]\) 内。\(n=20,\,S=0\) 时 Wilson 95\% 区间约为 \([0,\,0.16]\)---比 \([0,0]\) 诚实得多。

Clopper--Pearson 区间：基于精确二项分布，保证对每个 \(\theta\) 的覆盖率都 \emph{至少} \(1-\alpha\)（保守覆盖）。代价是区间偏宽。

似然比区间：收集所有不被似然比检验拒绝的参数值。有限样本行为通常介于 Wald 和 score 之间。

Bayesian Beta 后验区间：用 Beta 先验得出后验分位数。解释是``给定数据和先验，参数有 95\% 概率在里面''，和频率覆盖是两套逻辑。

每一种方法在保守性、平均长度、边界表现和解释上各有取舍。没有一个方法可以只凭``看起来合理''就保证拥有标称覆盖率。

\subsection{覆盖率不是一个笼统的数字}

``95\%''这个词下面掩盖了好几种不同的含义：

\begin{itemize}
\tightlist
\item 对每一个 \(\theta\) 都至少 95\% 的 uniform coverage（均匀覆盖）；
\item 仅在大样本极限下趋近 95\% 的 pointwise asymptotic coverage（逐点渐近覆盖）；
\item 对某个参数子空间或在某个条件事件下的 conditional coverage（条件覆盖）；
\item 经过多重比较校正后的 simultaneous coverage（同时覆盖），保证所有区间同时盖住各自参数的概率不低于 95\%。
\end{itemize}

只写``95\% 置信区间''而不说明是哪种覆盖、覆盖对象是谁、在什么条件下成立，会把一个精密的概念糊成一个空洞的标签。

\subsection{窄不等于可靠}

一个置信区间可以极窄，却不值得信任。窄区间的来源不只有``样本量大''和``噪声小''这两个正当理由：

\begin{itemize}
\tightlist
\item 模型错误地忽略了组内相关性，导致标准误被严重低估；
\item 在多个终点中选择最有利的那个报告，选择性报告把区间人为收窄；
\item 把超参数估计值当成已知常数，漏掉超参数的不确定性；
\item 数值软件在计算协方差时少传递了一层不确定性。
\end{itemize}

所以评价一个区间方法时，你需要问的问题都在看到具体端点之前：

\begin{itemize}
\tightlist
\item 覆盖率是精确的还是渐近的？有限样本下的实际覆盖率在什么水平？
\item 对哪些参数值有效？边界附近的表现怎样？
\item 平均长度和最坏情况长度分别如何？
\item 是否在数据看过之后才决定报告哪个参数、哪个区间？
\item 聚类、时间依赖和有限总体修正是否进入了标准误的计算？
\item 同时报告多个区间时，是否做了 family-wise 或 false coverage rate 的控制？
\end{itemize}

最后一个问题特别容易在日常分析中被跳过。如果你同时估计 20 个参数并分别给出 95\% 区间，即使每个区间单独都有标称覆盖，至少一个区间没盖住的概率可能远大于 5\%。这不是哪个区间出了问题，而是你在做 20 次抽奖，每次都觉得自己赢面很大，但总有一次会失手。

\subsection{从程序到这次区间}

回到开头那句话。``参数有 95\% 概率在 \([0.86,0.97]\) 里''---在频率学派框架内技术上不成立。但这不代表置信区间没有价值。一个诚实的 95\% 置信区间意味着：如果所有模型假设都站得住，在你能想象的所有重复实验中，使用这个程序的人有 95\% 的长期经验是区间盖住了参数。

这条保证不问你信不信先验，不要求参数服从某个分布。它只要求你把注意力从``这一次的数字''移开，先问清楚：产生这些数字的程序，在你说出``95\%''这两个字之前，有没有资格给出这个保证。
